\documentclass[11pt,a4paper]{article}
% =====================================================================
%  Shared preamble for the Part V lecture notes (lectures 19-22).
%
%  These four lectures sit outside Geron, so the notes ARE the primary
%  source and are examined on the same terms as the textbook parts.
%  Everything here is chosen to make that readable on paper: one column,
%  generous measure, and the same six-colour palette as the slides so a
%  student moving between the two is not re-learning what red means.
%
%  Build with pdflatex (twice, for the toc and the refs):
%      cd notes && latexmk -pdf lecture-19.tex
%  or  make -C notes
% =====================================================================
\usepackage[T1]{fontenc}
\usepackage[utf8]{inputenc}
\usepackage{newtxtext,newtxmath}      % Times-like: prints well, reads well
\usepackage[scaled=0.86]{inconsolata} % code, at a size that sits beside Times
\usepackage{microtype}
% newtx defines \Bbbk, \openbox, \proof and \endproof; amssymb and amsthm
% then redefine them and abort. Freeing the four names before those two
% packages load is the standard fix, and it costs nothing here: we use
% amsthm only for \newtheorem.
\let\Bbbk\relax
\usepackage{amsmath,amssymb}
\let\openbox\relax
\let\proof\relax
\let\endproof\relax
\usepackage{amsthm}
\usepackage{booktabs}
\usepackage{enumitem}
\usepackage{xcolor}
\usepackage{tcolorbox}
\usepackage{titlesec}
\usepackage{graphicx}
\usepackage[hidelinks]{hyperref}
\tcbuselibrary{skins,breakable}

% ---- the course palette, from assets/css/custom.css -------------------
\definecolor{aimlprimary}{HTML}{0B3D62}   % deep blue  - structure
\definecolor{aimlaccent} {HTML}{C0392B}   % brick red  - failures
\definecolor{aimlsuccess}{HTML}{14663A}   % green      - repairs
\definecolor{aimlmath}   {HTML}{6C3483}   % purple     - the derivation
\definecolor{aimlmuted}  {HTML}{4B5563}
\definecolor{aimlrule}   {HTML}{B0BCC7}
\definecolor{aimlsoft}   {HTML}{F4F7F9}

\usepackage[a4paper,margin=32mm,top=30mm,bottom=30mm]{geometry}
\setlength{\parskip}{0.45\baselineskip}
\setlength{\parindent}{0pt}
\linespread{1.06}

% ---- headings --------------------------------------------------------
\titleformat{\section}{\Large\bfseries\color{aimlprimary}}{\thesection}{0.7em}{}
\titleformat{\subsection}{\large\bfseries\color{aimlprimary}}{\thesubsection}{0.7em}{}
\titleformat{\subsubsection}{\bfseries\color{aimlmuted}}{}{0em}{}
\titlespacing*{\section}{0pt}{1.6\baselineskip}{0.5\baselineskip}
\titlespacing*{\subsection}{0pt}{1.2\baselineskip}{0.35\baselineskip}

% ---- boxes -----------------------------------------------------------
% One box per role, and the role is the colour. A student who has seen the
% slides already knows what each one means before reading a word of it.
\newtcolorbox{keybox}[1][]{
  breakable, enhanced, colback=aimlsoft, colframe=aimlprimary,
  boxrule=0.4pt, left=8pt, right=8pt, top=6pt, bottom=6pt,
  fonttitle=\bfseries\small, coltitle=white, colbacktitle=aimlprimary,
  attach boxed title to top left={yshift=-2mm, xshift=4mm},
  boxed title style={boxrule=0pt, arc=1pt}, #1}

\newtcolorbox{failbox}[1][]{
  breakable, enhanced, colback=white, colframe=aimlaccent,
  boxrule=0.9pt, left=8pt, right=8pt, top=6pt, bottom=6pt,
  fonttitle=\bfseries\small, coltitle=white, colbacktitle=aimlaccent,
  attach boxed title to top left={yshift=-2mm, xshift=4mm},
  boxed title style={boxrule=0pt, arc=1pt}, #1}

\newtcolorbox{derivbox}[1][]{
  breakable, enhanced, colback=white, colframe=aimlmath,
  boxrule=0.6pt, left=8pt, right=8pt, top=6pt, bottom=6pt,
  fonttitle=\bfseries\small, coltitle=white, colbacktitle=aimlmath,
  attach boxed title to top left={yshift=-2mm, xshift=4mm},
  boxed title style={boxrule=0pt, arc=1pt}, #1}

% An exercise. Deliberately the same shape as the notebook's `try` field:
% one change, and what should happen to the output.
\newtcolorbox{trybox}[1][]{
  breakable, enhanced, colback=white, colframe=aimlsuccess,
  boxrule=0.5pt, left=8pt, right=8pt, top=6pt, bottom=6pt,
  fonttitle=\bfseries\small, coltitle=white, colbacktitle=aimlsuccess,
  attach boxed title to top left={yshift=-2mm, xshift=4mm},
  boxed title style={boxrule=0pt, arc=1pt}, #1}

% ---- small semantics -------------------------------------------------
\newcommand{\scope}[1]{\textcolor{aimlmuted}{\small #1}}
\newcommand{\term}[1]{\textbf{#1}}
\newcommand{\code}[1]{\texttt{\small #1}}
\newcommand{\lecture}[1]{Lecture~#1}

\newtheorem{definition}{Definition}[section]
\newtheorem{proposition}[definition]{Proposition}

% ---- title block -----------------------------------------------------
% \notestitle{19}{Information retrieval: the lexical foundation}{SciFact (BEIR)}
\newcommand{\notestitle}[3]{%
  \thispagestyle{empty}
  \begin{flushleft}
    {\color{aimlmuted}\small\sffamily
     APPLICATIONS OF MACHINE LEARNING \quad$\cdot$\quad
     BSc Mathematics of Artificial Intelligence}\\[2mm]
    {\color{aimlrule}\rule{\textwidth}{0.8pt}}\\[4mm]
    {\color{aimlmuted}\large Lecture #1 \quad$\cdot$\quad Part V \quad$\cdot$\quad
     Extended notes}\\[3mm]
    {\Huge\bfseries\color{aimlprimary}\baselineskip=1.15\baselineskip #2\par}\vspace{4mm}
    {\color{aimlmuted}\large Dataset: #3}\\[3mm]
    {\color{aimlrule}\rule{\textwidth}{0.8pt}}
  \end{flushleft}
  \vspace{2mm}
  \begin{keybox}[title={Why these notes exist}]
    Lectures 19--22 sit outside G\'eron, the course's primary text. For those
    four lectures \emph{these notes are the primary source}, and they are
    examined on exactly the same terms as the textbook parts. Everything in
    the slide deck for this lecture is here, with the arguments written out at
    length rather than compressed onto a slide; the numbers are the ones the
    accompanying notebook prints, and every one of them is a measurement on
    the corpus named above rather than a figure quoted from a paper.
  \end{keybox}
  \vspace{1mm}
  {\small\tableofcontents}
  \vspace{2mm}
  \noindent{\color{aimlrule}\rule{\textwidth}{0.4pt}}
  \vspace{1mm}
}


\begin{document}
\notestitlefor{11}{II}{Training deep networks}{CIFAR-10}{Chapter 11}

\section{A deep stack that will not train}

You can build a network and you own its training loop. So build a deep one ---
twenty hidden layers on CIFAR-10 --- and train it exactly as Lecture 10 taught.

It does not train. The loss barely moves. Nothing raises, nothing is
\code{NaN}, the harness is correct and the data is fine, and a shallower
version of the same network trains without complaint.

Today we measure what goes wrong, derive why, and repair it --- in that order,
because the repair is unreadable until you have seen the measurement.

\subsection{The specification, and what it does not say}

3{,}072 inputs, 20 hidden layers of 100 units, a logistic activation on each, a
linear head of 10 outputs. Cross-entropy, Adam at 0.001, batches of 128, 20
epochs.

Read that again and note what it does not mention: \emph{how the weights start
out}. That is not something being hidden --- it is the ordinary case. Almost no
specification mentions it.

\begin{center}
\small
\begin{tabular}{llr}
\toprule
\textbf{Layer} & \textbf{Shape} & \textbf{Parameters} \\
\midrule
First & $3{,}072 \to 100$ & 307{,}300 \\
Each of the other 19 & $100 \to 100$ & 10{,}100 \\
Head & $100 \to 10$ & 1{,}010 \\
\midrule
Total & 21 weight matrices & 500{,}210 \\
\bottomrule
\end{tabular}
\end{center}

Comparable to Lecture 9's network, which reached 88\% on Fashion-MNIST.
Nothing here is under-parameterised.

A model that has learned nothing scores $\ln 10 = 2.3026$. Epoch 1 gives
2.3091 and epoch 20 gives 2.3036: the loss moved by 0.0055 in twenty epochs,
and the test accuracy is \textbf{10.0\%} --- exactly what a single
\code{return 3} would have bought.

\begin{center}
\small
\begin{tabular}{rrrr}
\toprule
\textbf{Hidden layers} & \textbf{Parameters} & \textbf{Final loss} &
\textbf{Test accuracy} \\
\midrule
 1 & 308{,}310 & 1.0259 & 41.2\% \\
 2 & 318{,}410 & 1.0503 & 39.3\% \\
 5 & 348{,}710 & 1.6925 & 21.0\% \\
10 & 399{,}210 & 2.3035 & 10.0\% \\
20 & 500{,}210 & 2.3036 & 10.0\% \\
\bottomrule
\end{tabular}
\end{center}

Adding layers made it worse and past a point made it useless --- which is not
what capacity is supposed to do. And the parameter count barely moves after the
first layer, because the 3{,}072-input layer dominates, so this is not a
capacity effect.

The loop is correct, the data is correct, the metric is correct, and the depth
is the variable. Which is where Lecture 10's real purchase pays.

\section{Instrumenting it}

Backpropagation computes 500{,}210 partial derivatives on every batch. Under
\code{MLPClassifier.fit()} every one was discarded; they are all still here in
\code{p.grad}. Two probes, about ten lines each --- without them the diagnosis
is ``it does not train'' and there is no next step.

\subsection{Probe 1, and why it was wrong}

Running the modules by hand and recording each logistic's output:

\begin{center}
\small
\begin{tabular}{rrrr}
\toprule
\textbf{Layer} & \textbf{Mean} & \textbf{Standard deviation} &
\textbf{Saturated} \\
\midrule
 1 & near 0.50 & 0.1311 & 0.000 \\
 2 & near 0.50 & 0.0707 & 0.000 \\
 5 & near 0.50 & 0.0773 & 0.000 \\
10 & near 0.50 & 0.0731 & 0.000 \\
20 & near 0.50 & 0.0710 & 0.000 \\
\bottomrule
\end{tabular}
\end{center}

Mean near a half, spread near 0.071, flat for twenty layers, nothing saturated.
The forward signal looks fine, and the textbook's first explanation does not
apply.

\begin{failbox}[title={The probe measured the wrong thing}]
\code{h.std()} spans the \emph{whole tensor}, so it is dominated by the spread
of the layer's random biases across units --- which does not change with depth.
It would read 0.07 in a network that had stopped computing entirely.

What carries information is the spread \emph{down the batch}: how much a unit
moves when the input does.

\begin{center}
\small
\begin{tabular}{lrrrr}
\toprule
\textbf{Hidden layer} & 1 & 5 & 10 & 20 \\
\midrule
\code{h.std()} --- what we plotted & 0.131 & 0.071 & 0.071 & 0.071 \\
\code{h.std(dim=0).mean()} --- the signal &
   $1.3\mathrm{e}{-1}$ & $5.9\mathrm{e}{-5}$ & $3.0\mathrm{e}{-9}$ &
   $2.4\mathrm{e}{-17}$ \\
\bottomrule
\end{tabular}
\end{center}

Sixteen orders of magnitude. The forward signal \emph{does} die, by a factor of
0.149 per layer, and by layer 20 the network is no longer a function of its
input.

The saturation column was no better. It asks whether
$|h - \tfrac12| > 0.45$ in a band whose standard deviation is 0.071 --- a
threshold 6.3 standard deviations away, and 3.4 even at layer 1. Nothing could
ever have exceeded it. ``0.000 at every depth'' was arithmetic, not evidence.
\end{failbox}

\subsection{Probe 2 --- the backward pass}

Eight batches rather than one, because a single batch of 128 is a noisy
estimate of anything and a course that says so should not then quote one.

\begin{center}
\small
\begin{tabular}{lr}
\toprule
\textbf{Layer} & $\lVert\partial L/\partial W\rVert$ \\
\midrule
1 --- nearest the input & $4.567\mathrm{e}{-17}$ \\
5   & $9.296\mathrm{e}{-15}$ \\
10  & $1.845\mathrm{e}{-10}$ \\
15  & $3.042\mathrm{e}{-06}$ \\
20  & $6.608\mathrm{e}{-02}$ \\
head & $4.866\mathrm{e}{-01}$ \\
\bottomrule
\end{tabular}
\end{center}

Layer 20 receives $1.45\times10^{15}$ times what layer 1 does.

\begin{keybox}[title={The shape of that line is the finding}]
It is \emph{straight} on a logarithmic axis, and a straight line on a log axis
is a geometric sequence. The attenuation is not one bad layer or a fault at the
input: it is the same factor applied at every layer,
\[ \frac{\lVert g_\ell\rVert}{\lVert g_{\ell+1}\rVert} \approx \rho
   \qquad\Longrightarrow\qquad
   \frac{\lVert g_1\rVert}{\lVert g_{20}\rVert} \approx \rho^{19}. \]
One number characterises the whole network, and here it is
\[ \rho = 0.1593, \qquad \rho^{19} = 6.91\times10^{-16}. \]
\end{keybox}

Does training rescue it? We logged the per-layer norms every epoch, so this is
a lookup rather than a guess: layer 1's gradient was
$6.660\mathrm{e}{-17}$ at epoch 1 and $4.610\mathrm{e}{-17}$ at epoch 20.
Twenty epochs of Adam did not move it into a range a learning rate of 0.001
could use. And measuring
$\lVert W_{\text{after}} - W_{\text{before}}\rVert / \lVert W_{\text{before}}
\rVert$, layer 1 moved 0.0000, layer 10 moved 0.0001 and layer 20 moved 0.5890. Only
the top six layers learned anything, and they had nothing useful beneath
them.

\begin{failbox}[title={Two replies that do not work}]
\emph{Adam will compensate.} It rescales, so the direction is what matters ---
and at $10^{-17}$ the direction is dominated by batch-to-batch noise. Worse,
Adam's $\varepsilon$, default $10^{-8}$, is \emph{larger} than the gradients at
the bottom of the stack. An optimiser can compensate for a badly scaled
gradient; it cannot manufacture information the backward pass did not deliver.

\emph{Raise the learning rate.} The head's gradient is
$4.87\mathrm{e}{-01}$ and layer 1's is $4.57\mathrm{e}{-17}$, so any rate large
enough to move layer 1 is $10^{16}$ times too large for the head. A single
scalar cannot serve a quantity spanning sixteen orders of magnitude. The repair
has to change the \emph{spread}, not the scale.
\end{failbox}

The chapter names two failures, and they are the same mechanism at different
values of one number. Vanishing is $\rho < 1$: the loss does not move, the
layers nearest the data suffer, and it is silent. Exploding is $\rho > 1$:
\code{NaN} or a jumping loss, all layers at once, and it is loud. We have the
quiet one --- and one derivation covers both.

\section{Derivation: variance propagation through layers}

Why 0.1593? Not ``why is it small'' --- why is it \emph{that number}, and what
would it have to be instead?

Drop the activation and look at the linear part:
$z_i = \sum_{j=1}^{n_{\text{in}}} w_{ij}a_j$. Three assumptions, stated so they
can be checked: the weights are drawn independently with mean zero; they are
independent of the inputs (true at initialisation, approximately true
afterwards); and the inputs have mean zero and are uncorrelated.

All three are idealisations, which is fine --- the point of an idealisation is
that you can compute with it and then check the answer. Assumption 3 is the
weakest at the input layer, since Lecture 9 divided the pixels by 255 so they
arrive in $[0,1]$ with a mean near 0.5.

\begin{derivbox}[title={The one recursion}]
Each term $w_{ij}a_j$ has mean zero and the terms are uncorrelated, so the
variances add:
\[ \operatorname{Var}(z)
   = n_{\text{in}}\cdot\operatorname{Var}(w)\cdot\operatorname{Var}(a). \]
Everything in this lecture is a consequence of that line. Note its shape: the
input variance is \emph{multiplied} by a factor, which is what makes depth
dangerous.
\end{derivbox}

Check it before believing it --- one layer, fan-in 100, twenty thousand random
inputs of unit variance:

\begin{center}
\small
\begin{tabular}{lrr}
\toprule
$\operatorname{Var}(w)$ & \textbf{Predicted $\operatorname{Var}(z)$} &
\textbf{Measured} \\
\midrule
0.001 & 0.1000 & 0.1002 \\
0.050 & 5.0000 & 5.0050 \\
\bottomrule
\end{tabular}
\end{center}

\subsection{Two requirements, and they conflict}

The forward pass wants the signal to arrive at layer 20 with the spread it had
at layer 1:
\[ \operatorname{Var}(z) = \operatorname{Var}(a)
   \iff n_{\text{in}}\operatorname{Var}(w) = 1
   \iff \operatorname{Var}(w) = \frac{1}{n_{\text{in}}}. \]

The backward pass runs the same matrix \emph{transposed}. Using Lecture 10's
adjoint notation, $\bar a = W^{\mathsf T}\bar z$, so identical algebra with one
index swapped:
\[ \operatorname{Var}(\bar a)
   = n_{\text{out}}\cdot\operatorname{Var}(w)\cdot\operatorname{Var}(\bar z)
   \qquad\Longrightarrow\qquad
   \operatorname{Var}(w) = \frac{1}{n_{\text{out}}}. \]

The index swaps because a row of $W$ and a column of $W$ are different sums:
forward, each output unit sums over $n_{\text{in}}$ terms; backward, each input
unit collects from $n_{\text{out}}$. That is exactly Lecture 10's fan-in /
fan-out bookkeeping --- a node's adjoint is the sum over its outgoing edges,
and here we are counting how many there are.

\begin{center}
\small
\begin{tabular}{lrrl}
\toprule
\textbf{Layer} & $n_{\text{in}}$ & $n_{\text{out}}$ &
\textbf{Can both be satisfied?} \\
\midrule
First & 3{,}072 & 100 & no --- they differ by a factor of 30 \\
Each hidden & 100 & 100 & yes, exactly \\
Head & 100 & 10 & no \\
\bottomrule
\end{tabular}
\end{center}

Every network whose layers change width --- which is every network --- has to
give something up. \term{Glorot} does not choose, but takes the harmonic mean:
\[ \operatorname{Var}(w) = \frac{2}{n_{\text{in}} + n_{\text{out}}}, \]
a small error in both directions rather than a large error in one. Harmonic
because the quantities being averaged are reciprocals of the fans, and on a
square layer it is not a compromise at all --- it reduces to $1/n$. On our
hidden layers that is $\operatorname{Var}(w) = 0.0100$, a standard deviation of
0.1000, against the default's 0.0577.

\subsection{Putting the activation back}

Between two linear maps sits $\varphi$, and the backward pass multiplies by its
derivative pointwise, so
\[ \rho = \sqrt{\,n_{\text{out}}\cdot\operatorname{Var}(w)\cdot
                 \mathbb{E}\bigl[\varphi'(z)^2\bigr]\,}. \]
Two factors we control --- the shape of the matrix and the variance we
initialise with --- and one belonging to the activation.

\begin{failbox}[title={The logistic's ceiling}]
$\sigma'(z) = \sigma(z)(1-\sigma(z)) \le \tfrac14$ for \emph{every} $z$, with
equality only at $z = 0$. Not ``usually small'': bounded above by a quarter
everywhere.

So the logistic multiplies $\rho$ by at most 0.25 before any weight has been
chosen. Even a perfectly initialised logistic stack cannot do better than
$(1/4)^{19}$.
\end{failbox}

ReLU discards exactly half the variance instead: for symmetric zero-mean $z$,
$\mathbb{E}[\operatorname{ReLU}'(z)^2] = \tfrac12$ --- half the time the
derivative is 1, half the time 0. A clean factor of two, and a factor of two
can be paid for. \term{He} pays it by doubling the weight variance,
$\operatorname{Var}(w) = 2/n_{\text{in}}$, giving
$\rho = \sqrt{n_{\text{out}}/n_{\text{in}}}$ --- exactly 1 when the fans are
equal. He fixes the \emph{forward} condition, so it uses the fan-in; off square
layers, Glorot's compromise is still the honest choice.

\begin{center}
\small
\begin{tabular}{lrrr}
\toprule
\textbf{Scheme} & $\operatorname{Var}(w)$ &
$\mathbb{E}[\varphi'^2]$ & \textbf{Predicted $\rho$} \\
\midrule
default, logistic & 0.00333 & 0.0625 & 0.1443 \\
Glorot, logistic  & 0.01000 & 0.0625 & 0.2500 \\
Glorot, ReLU      & 0.01000 & 0.5000 & 0.7071 \\
He, ReLU          & 0.02000 & 0.5000 & 1.0000 \\
\bottomrule
\end{tabular}
\end{center}

Nothing in that table came from running anything. It is the fan-in, the
fan-out, and one expectation --- and the first row is the 0.1593 we measured.

\begin{keybox}[title={And then it compounds}]
\[ \frac{\lVert g_1\rVert}{\lVert g_L\rVert} = \rho^{\,L-1}
   \qquad \text{$L$ layers, $L-1$ steps between them.} \]

\begin{center}
\small
\begin{tabular}{lll}
\toprule
$\rho$ & \textbf{over 19 steps} & \textbf{what you see} \\
\midrule
0.144 & $10^{-16}$ & the loss does not move \\
0.707 & $10^{-3}$  & trains, slowly, badly \\
1.000 & 1          & every layer learns \\
1.4 & $10^{3}$   & \code{NaN} \\
\bottomrule
\end{tabular}
\end{center}
\emph{There is no safe value of $\rho$ except one.} That is what makes deep
networks a separate subject rather than a bigger version of shallow ones.
\end{keybox}

\section{Seven repairs, measured one at a time}

\subsection{Initialisation, and why it is not enough}

\begin{center}
\small
\begin{tabular}{lrrr}
\toprule
& $\rho$ & \textbf{End to end} & \textbf{Test accuracy} \\
\midrule
Default, logistic & 0.1593 & $1.4\mathrm{e}{+15}$ & 10.0\% \\
Glorot, logistic  & 0.25 & $2.2\mathrm{e}{+11}$ & 10.0\% \\
\bottomrule
\end{tabular}
\end{center}

Nearly four orders of magnitude better, and still eleven orders of attenuation
--- which is hopeless. The measured $\rho$ is now 0.25 --- which is \emph{exactly} the logistic's own
ceiling, with the weight scale perfectly neutralised. \emph{You cannot initialise your way out of the activation
function.}

\subsection{The activation, and normalisation}

Every requirement on $\varphi$ comes from the derivation: $\varphi'$ should
reach 1 somewhere and not be bounded below it; $\varphi$ should not saturate on
both sides; and the output should be roughly centred so assumption 3 holds at
the next layer. ReLU satisfies the first two and fails the third, and is still
the right default --- though it has its own failure mode, a unit whose
pre-activation is negative for every input has gradient exactly zero forever.
Leaky ReLU, ELU and SELU all repair that the same way, by giving the negative
side a non-zero slope.

\begin{center}
\small
\begin{tabular}{llr}
\toprule
\textbf{Activation} & \textbf{Initialisation} & \textbf{Test accuracy} \\
\midrule
logistic & Glorot & 40.4\% \\
tanh & Glorot & 38.5\% \\
ReLU & He & 36.8\% \\
Leaky ReLU & He & 37.1\% \\
ELU & He & 43.2\% \\
SELU & Glorot & 41.0\% \\
\bottomrule
\end{tabular}
\end{center}

Spread 6.4 points --- and the logistic is \emph{not last}. All six carry batch
normalisation, which is why it survives: normalisation puts the pre-activations
back where its derivative is largest, at every step. Without normalisation,
moving off the logistic was worth 31.1 points. With it, the choice of
activation is worth 6.4.

Batch normalisation standardises each layer's inputs and then learns what scale
they should have had,
\[ \hat z = \frac{z - \mu_B}{\sqrt{\sigma_B^2 + \varepsilon}},
   \qquad y = \gamma\hat z + \beta, \]
where the second step is what stops it being a straitjacket --- the network can
undo the normalisation if it needs to.

\begin{center}
\small
\begin{tabular}{lrrr}
\toprule
\textbf{On the ReLU + He network} & \textbf{Test accuracy} &
\textbf{Parameters} & \textbf{Epoch cost} \\
\midrule
No normalisation & 41.1\% & 500{,}210 & baseline \\
Batch normalisation & 36.8\% & 504{,}210 & 2.16 times the epoch \\
Layer normalisation & 41.6\% & 504{,}210 & cheaper, and still not free \\
\bottomrule
\end{tabular}
\end{center}

\begin{keybox}[title={Why the repair everyone reaches for first did not pay}]
Batch normalisation applied to the \emph{broken} network takes it from 10.0\%
to 37.9\%. Applied to the \emph{already repaired} network it takes 41.1\% down
to 36.8\%, and costs $2.16\times$ the epoch.

Normalisation exists to hold $\rho$ near 1 when something else has pushed it
away. He initialisation had already put it there, so there was nothing left to
correct --- and what remained was the cost. On a wider network or a longer run
the trade goes the other way. The point is that you can tell which case you are
in, because you measured both.
\end{keybox}

\begin{failbox}[title={The bill batch normalisation carries}]
Its statistics depend on the other images in the batch, so at training time the
prediction for one image depends on which images happened to be batched with
it, while at evaluation it uses running averages. Those are two different
functions.

That is Lecture 10's \code{model.eval()} failure again --- and now much harder
to catch. Dropout wobbles between calls, so running the evaluation twice
exposes it; batch normalisation in training mode is deterministic for a fixed
batch, so the same trick finds nothing.
\end{failbox}

\subsection{Clipping, optimisers and schedules}

Gradient clipping rescales the whole gradient vector when its norm exceeds a
threshold, preserving the direction and bounding the length. It is a defence
against $\rho > 1$, and we have $\rho < 1$: applied alone to the failing
network it buys nothing at all. On the repaired network it took 36.8\% to
39.6\%, but the median step norm with He is 3.19 and the threshold was 1.0 ---
so it acted as a step-size limiter rather than a safety net. \emph{A useful
effect obtained for the wrong reason}; if you want a smaller step, set a
smaller learning rate and say so.

Comparing optimisers on the broken network would have measured nothing. On the
repaired one: SGD 27.5\%, momentum 33.6\%, Nesterov 35.2\%, RMSProp 36.8\%,
Adam 36.8\%, AdamW 36.1\%. The spread between best and worst is far smaller
than the effect of the initialisation --- order your effort by the measurement,
not by the novelty.

Schedules need to know how long the run is, which is why changing the epoch
count silently changes the schedule. Constant 39.6\%, cosine 38.8\%, 1-cycle
41.5\% --- on one seed each, so read that as a shape rather than a ranking.

\section{The ablation}

Each row adds one change to the row above. Five seeds a row.

\begin{center}
\small
\begin{tabular}{lrrr}
\toprule
\textbf{Configuration} & \textbf{Test accuracy} & \textbf{sd} &
\textbf{Change} \\
\midrule
the failing network, unchanged & 10.0\% & $\pm$0.00 & --- \\
+ Glorot initialisation & 10.0\% & $\pm$0.00 & $+0.0$ (noise) \\
+ ReLU and He initialisation & 41.3\% & $\pm$0.35 & $\mathbf{+31.3}$ \\
+ batch normalisation & 36.5\% & $\pm$0.64 & $-4.7$ \\
+ gradient clipping & 40.8\% & $\pm$0.76 & $+4.3$ \\
+ a 1-cycle schedule & 41.3\% & $\pm$2.79 & $+0.5$ (noise) \\
+ dropout 0.1 & 34.0\% & $\pm$1.33 & $-7.3$ \\
\bottomrule
\end{tabular}
\end{center}

The third column is the deliverable, and the fourth is what makes it readable.
Typical spread is 0.84 points; one step is worth $+31.3$, forty times that.
Three more clear it. The 1-cycle schedule does not --- $+0.5$ against
$\pm2.79$ --- \emph{and the single-seed version of this table called it the
winner}.

Applied alone to the failing network: Glorot 10.0\%, ReLU with He 41.1\%, batch
normalisation 37.9\%, layer normalisation 10.0\%, clipping 10.0\%, the schedule
10.0\%, dropout 10.0\%. Four of the seven leave the network exactly where it
was --- which is not a disappointment but what tells you what each one is for.
Clipping, the schedule and dropout answer problems this network does not have,
and Glorot buys eleven orders of magnitude of gradient and is still not enough.

\begin{keybox}[title={The answer to ``should I add batch normalisation?''}]
One change is worth more than the whole ladder. The question is malformed until
you know which failure you have --- ours was $\rho \ne 1$, and only the repairs
that address $\rho$ paid.
\end{keybox}

\subsection{The silent-failure catalogue so far}

\begin{center}
\small
\begin{tabular}{llr}
\toprule
\textbf{Failure} & \textbf{Diagnosed in} & \textbf{Measured cost} \\
\midrule
Scaler fitted before the split & Lecture 2 & nothing measurable \\
Accuracy under class imbalance & Lecture 3 & a useless model above 90\% \\
Missing \code{optimizer.zero\_grad()} & Lecture 10 & 76.7 points \\
Metric averaged per batch & Lecture 10 & 0.382 points \\
A deep stack on default initialisation & today & 31.3 points \\
Glorot's constant used with ReLU & today & 1.5 points \\
\bottomrule
\end{tabular}
\end{center}

Every one ran to completion and printed a plausible number. \emph{The last two
are the first in which no line of code is wrong.}

\section{The notebook}

\begin{trybox}[title={What to change}]
Set the initialisation back to the default and re-run the gradient probe. The
straight line on a log axis comes back, and its slope is the per-layer factor
the derivation predicts --- so you will have closed the loop from measurement
to theory to measurement yourself.
\end{trybox}

\section{Exercises}

\begin{enumerate}[leftmargin=1.6em,itemsep=4pt]
  \item One layer multiplies the standard deviation of the backward signal by
        $\rho = \sqrt{n_{\text{out}}\operatorname{Var}(w)
        \mathbb{E}[\varphi'^2]}$. State what $L$ layers do to it, and why only
        $\rho = 1$ survives depth. \hfill \scope{[5]}
  \item Preserving the forward signal asks for
        $\operatorname{Var}(w) = 1/n_{\text{in}}$ and preserving the gradient
        asks for $1/n_{\text{out}}$. State when the two agree, and what Glorot
        does when they do not. \hfill \scope{[4]}
  \item He initialisation doubles the weight variance for ReLU. Derive the
        factor of two in one line. \hfill \scope{[4]}
  \item A gradient profile that is a straight line on a log axis tells you
        something a curved one does not. State what, and why it identifies the
        cause. \hfill \scope{[4]}
  \item Gradient clipping is applied to a network whose gradients are fifteen
        orders of magnitude too small. Predict the effect, and say what the
        attempt reveals about the diagnosis. \hfill \scope{[4]}
\end{enumerate}

\section*{Summary}
\addcontentsline{toc}{section}{Summary}

A twenty-layer stack sat at chance while nothing raised. Instrumenting it ---
which is what owning the loop bought --- showed the gradient shrinking by a
constant factor per layer, a straight line on a log axis, so one number
characterises the whole network: $\rho = 0.1593$, compounding to
$1.45\times10^{15}$ over nineteen steps.

The first probe was wrong, and the notes keep it. \code{h.std()} spans the
whole tensor and reads 0.07 in a network that has stopped computing; the spread
down the batch falls by sixteen orders of magnitude. And a saturation threshold
6.3 standard deviations from the mean returned 0.000 at every depth, which was
arithmetic rather than evidence.

One recursion explains it:
$\operatorname{Var}(z) = n\operatorname{Var}(w)\operatorname{Var}(a)$, with the
fan-in forwards and the fan-out backwards. The two are the same condition only
on a square layer, Glorot takes their harmonic mean, and He doubles it because
ReLU discards half the variance. The predicted $\rho$ for the four schemes
comes from the shapes alone, and the first row is the number we measured.

Seven repairs, each measured alone and then in a ladder over five seeds: the
step worth $+31.3$ points is ReLU with He, against a typical spread of 0.84.
Batch normalisation \emph{cost} 4.7 points on a network that no longer needed
it while doubling the epoch, and the 1-cycle schedule's $+0.5$ against
$\pm2.79$ was called a winner by the single-seed version of the same table.

\end{document}
